\documentclass[../general/general]{subfiles}

\begin{document}

\chapter{Лекция №2: Термодинамические следствия микроканонического ансамбля и вывод канонического распределения Гиббса}
\noindent\textit{А.\,М.~Белемук \hfill 20 февраля 2021}
\vspace{1em}
\ifSubfilesClassLoaded{\tableofcontents\newpage}{}

% --- SECTION 1: Связь микроканонического распределения с термодинамикой ---
\section{Связь микроканонического распределения с термодинамикой}

На прошлой лекции мы вычислили энтропию классического идеального газа в рамках микроканонического ансамбля:
\begin{equation}
S(E, V, N) = N \ln \left[ \frac{eV}{N(2\pi\hbar)^3} \left( \frac{4\pi m e E}{3N} \right)^{3/2} \right]
\end{equation}

Теперь давайте установим связь этого выражения с классической термодинамикой. Для этого нам понадобится фундаментальное термодинамическое соотношение, которое записывается для дифференциала энтропии $dS$:

\begin{equation}
dS = \frac{1}{T} dE + \frac{P}{T} dV - \frac{\mu}{T} dN
\end{equation}

Из этого соотношения следует физический смысл частных производных энтропии по соответствующим переменным. В частности, коэффициенты перед дифференциалами имеют следующий вид:

\begin{equation}
\left( \frac{\partial S}{\partial E} \right)_{V, N} = \frac{1}{T}, \quad \left( \frac{\partial S}{\partial V} \right)_{E, N} = \frac{P}{T}, \quad \left( \frac{\partial S}{\partial N} \right)_{E, V} = -\frac{\mu}{T}
\end{equation}

\subsection{Определение абсолютной температуры и внутренняя энергия}

Начнем с первой производной — по энергии. Это позволит нам определить температуру системы. Напомню, что такое определение температуры через производную энтропии по энергии можно обосновать несколькими способами. Например, в задании есть задача показать, что когда две системы приходят в тепловой контакт, то именно из условия максимальности энтропии суммарной системы следует равенство этих частных производных для каждой из систем. А равенство производных — это и есть условие теплового равновесия, то есть равенство температур.

Нам остается продемонстрировать, что введенная таким образом температура действительно совпадает с абсолютной термодинамической температурой. Сделаем это на примере нашего идеального газа. Возьмем производную от выражения для энтропии, полученного ранее:

\begin{equation}
\frac{1}{T} = \frac{\partial}{\partial E} \left( N \ln \left[ E^{3/2} \cdot \dots \right] + f(V, N) \right)
\end{equation}

Для удобства дифференцирования перепишем выражение для энтропии, выделив зависимость от энергии:
\begin{equation}
S = N \ln E^{3/2} + \text{функция, не зависящая от } E
\end{equation}

Дифференцируя по $E$, получаем:
\begin{equation}
\frac{1}{T} = N \cdot \frac{3}{2} \cdot \frac{1}{E} = \frac{3N}{2E}
\end{equation}

Отсюда мы немедленно находим выражение для внутренней энергии:
\begin{equation}
E = \frac{3}{2} NT
\end{equation}

Это всем нам знакомое выражение для внутренней энергии идеального одноатомного классического газа. Таким образом, мы подтвердили, что температура $T$, входящая в наше статистическое определение, — это та самая термодинамическая температура.

\subsection{Давление и уравнение состояния}

Теперь перейдем ко второму результату. Найдем давление газа, используя производную энтропии по объему при постоянных энергии и числе частиц:

\begin{equation}
\frac{P}{T} = \left( \frac{\partial S}{\partial V} \right)_{E, N}
\end{equation}

Снова обратимся к нашему выражению для $S$. Зависимость от объема там аддитивная под логарифмом: $S = N \ln V + \dots$. Дифференцируя по объему, получаем:

\begin{equation}
\frac{P}{T} = \frac{\partial}{\partial V} (N \ln V + \dots) = \frac{N}{V}
\end{equation}

Следовательно:
\begin{equation}
P = \frac{N}{V} T = nT
\end{equation}

Мы получили уравнение состояния идеального газа, где $n = N/V$ — концентрация частиц. Как мы видим, в статистической физике давление также определяется через производную энтропии.

\subsection{Химический потенциал}

Наконец, рассмотрим химический потенциал $\mu$. Он показывает, насколько изменяется энтропия при добавлении одной частицы в систему при фиксированных энергии и объеме. Вычислим производную по $N$:

\begin{equation}
-\frac{\mu}{T} = \left( \frac{\partial S}{\partial N} \right)_{E, V}
\end{equation}

Дифференцирование по $N$ несколько сложнее, так как $N$ входит и как множитель перед логарифмом, и в знаменатель под логарифмом. Если проделать эту операцию (включая дифференцирование слагаемых, возникших из формулы Стирлинга для $N!$), мы получим следующее выражение для химического потенциала классического газа:

\begin{equation}
\mu = -T \ln \left[ \frac{V}{N} \left( \frac{2\pi m T}{h^2} \right)^{3/2} \right]
\end{equation}

Здесь мы уже подставили вместо энергии ее выражение через температуру. Это выражение можно переписать более компактно, если ввести понятие удельного объема $v = V/N$ (объем, приходящийся на одну частицу) и некоторую величину с размерностью объема, которую мы позже назовем квантовым объемом.

Важно отметить, что для классического идеального газа под логарифмом стоит величина, которая в классическом пределе много больше единицы. Поэтому химический потенциал классического газа всегда отрицателен: $\mu < 0$. К более детальному обсуждению этого вопроса мы вернемся чуть позже.


% --- SECTION 3: Квантовый объем и границы применимости классического описания ---
\section{Квантовый объем и границы применимости классического описания}

Теперь давайте более подробно разберем выражение, которое у нас получилось для химического потенциала, и поймем его физический смысл. Мы записали химический потенциал в следующем виде:

\begin{equation}
\mu(T, V, N) = -T \ln \left[ \frac{V}{N} \left( \frac{2\pi m T}{h^2} \right)^{3/2} \right]
\end{equation}

Здесь под логарифмом стоит величина, которая обязана быть безразмерной. Давайте выделим в ней структурные части. Отношение $V/N$ мы обозначим как $v$ — это удельный объем, то есть объем фазового пространства, приходящийся на одну частицу. 

Вторая часть выражения под логарифмом имеет размерность обратного объема. Мы можем ввести величину $V_Q$, которую назовем квантовым объемом:

\begin{equation}
\frac{1}{V_Q} = \left( \frac{2\pi m T}{h^2} \right)^{3/2} = \left( \frac{\sqrt{2\pi m T}}{h} \right)^{3}
\end{equation}

Тогда химический потенциал переписывается очень изящно:

\begin{equation}
\mu = -T \ln \left( \frac{v}{V_Q} \right)
\end{equation}

\subsection{Тепловая длина волны де Бройля}

Чтобы понять происхождение этого «квантового объема», нужно вспомнить основы квантовой механики. Каждой частице с импульсом $p$ соответствует волна де Бройля с длиной волны $\lambda = h/p$. В статистической системе при температуре $T$ частицы обладают некоторым характерным тепловым импульсом $p_T$.

Характерный тепловой импульс частицы по порядку величины можно оценить из соотношения для кинетической энергии: $p^2/2m \sim T$ (в единицах константы Больцмана). Точнее, для статистических расчетов удобно ввести тепловой импульс как $p_T = \sqrt{2\pi m T}$. Тогда мы получаем определение тепловой длины волны де Бройля:

\begin{equation}
\lambda_T = \frac{h}{p_T} = \frac{h}{\sqrt{2\pi m T}}
\end{equation}

Легко видеть, что введенный нами ранее квантовый объем $V_Q$ — это не что иное, как куб тепловой длины волны де Бройля:

\begin{equation}
V_Q = \lambda_T^3
\end{equation}

% Место под график/рисунок
% Комментарий: Лектор рисует куб со стороной \lambda_T, иллюстрируя масштаб квантовой неопределенности частицы.
% Таймкод: [19:50]

\subsection{Условие классичности и знак химического потенциала}

Теперь мы можем сформулировать критерий, при котором наше классическое рассмотрение (газ Больцмана) вообще справедливо. Классическое описание работает тогда, когда среднее расстояние между частицами много больше их квантового размера, определяемого тепловой длиной волны.

Если удельный объем $v = V/N$ (объем «ячейки» на одну частицу) много больше квантового объема $V_Q$:
\begin{equation}
v \gg V_Q \quad \text{или} \quad \frac{V}{N} \gg \lambda_T^3
\end{equation}
то такой режим называется классическим. В этом случае аргумент логарифма в выражении для химического потенциала велик, сам логарифм положителен, и из-за знака «минус» перед формулой мы получаем, что для классического идеального газа:
\begin{equation}
\mu < 0
\end{equation}

Химический потенциал классического газа всегда отрицателен. Это важное следствие, к которому мы еще не раз будем возвращаться.

\subsection{Температура вырождения}

Давайте посмотрим, что происходит при понижении температуры. Глядя на формулу для квантового объема:
\begin{equation}
V_Q(T) = \frac{h^3}{(2\pi m T)^{3/2}}
\end{equation}
мы видим, что при $T \to 0$ квантовый объем $V_Q$ стремится к бесконечности. Это означает, что при достаточно низких температурах условие классичности $v \gg V_Q$ неизбежно нарушится.

Когда квантовый объем становится соизмеримым с удельным объемом:
\begin{equation}
V_Q(T_0) \sim \frac{V}{N}
\end{equation}
газ перестает быть классическим. Частицы начинают «чувствовать» квантовую природу друг друга, их волновые пакеты перекрываются. Эта характерная температура $T_0$ называется температурой вырождения. 

При температурах $T \lesssim T_0$ мы переходим в квантовый режим, где классическая статистика Больцмана более не применима, и нужно использовать статистику Бозе-Эйнштейна или Ферми-Дирака (в зависимости от спина частиц). Там возникают такие явления, как Бозе-конденсация и другие квантовые эффекты, которые мы будем изучать позже. Но в рамках данной лекции мы работаем в классическом приближении, полагая $T \gg T_0$.


% --- SECTION 4: Каноническое распределение (распределение Гиббса) ---
\section{Каноническое распределение (распределение Гиббса)}

Мы с вами рассмотрели микроканонический ансамбль для идеального газа и провели вычисления термодинамических величин. Для этого нам потребовалось вычислить фазовый интеграл $\Gamma(E)$, затем найти энтропию $S(E)$ и уже из неё получить температуру $T$, давление $P$ и химический потенциал $\mu$. Как вы видели, основной трудностью в этой схеме является именно вычисление фазовых интегралов. Даже для идеального газа это потребовало определенных усилий, а для более сложных систем это может стать непреодолимой преградой.

Однако статистическая физика не стояла на месте, она развивалась. Был придуман и введен метод, который позволяет изящно обойти эти трудности и находить все термодинамические функции гораздо проще. По сути, этот метод предложил Гиббс, и тема нашей второй лекции — \textbf{Каноническое распределение} или \textbf{Распределение Гиббса}.

\subsection{Постановка задачи: тело в среде}

Теперь мы переходим к рассмотрению систем, которые более приближены к реальности. Раньше мы рассуждали на языке изолированных систем, которые практически не взаимодействуют с окружением. В реальности же любая система (тело) находится в некоторой внешней среде.

Представим нашу исследуемую систему как «тело» (обозначим индексом 1), для которого заданы объем $V$ и число частиц $N$. Это те параметры, которые мы контролируем. Тело находится в тепловом контакте с гораздо более крупной системой — «средой», «резервуаром» или «термостатом» (обозначим индексом 2).

% Место под рисунок
% Комментарий: Лектор рисует небольшое тело (прямоугольник с цифрой 1) внутри огромной замкнутой области (среда с цифрой 2). Указано, что полная система (1+2) изолирована.
% Таймкод: [20:09]

Этот термостат задает и фиксирует для нашего тела температуру $T$. Тело не изолировано от среды, оно обменивается с ней энергией. 

\subsection{Hamilton-функция и условие слабого взаимодействия}

Рассмотрим полную систему, состоящую из тела и среды вместе. Эту полную систему $(1+2)$ мы будем считать изолированной. Её полная энергия $E$ постоянна (точнее, лежит в узком интервале от $E$ до $E + \Delta E$).

Hamilton-функция всей системы $H(p, q)$ зависит от координат и импульсов частиц тела $(p_1, q_1)$ и частиц среды $(p_2, q_2)$. В общем случае она записывается как:
\begin{equation}
H(p, q) = H_1(p_1, q_1) + H_2(p_2, q_2) + V_{12}(p, q)
\end{equation}
где $V_{12}$ — энергия взаимодействия между подсистемами.

Здесь есть важный нюанс. Мы предполагаем наличие \textbf{слабого взаимодействия} между телом и средой. Оно считается настолько маленьким, что практически не сдвигает уровни энергии Hamilton-функции $H_1$, но при этом оно \textit{обязательно} должно присутствовать. Именно взаимодействие $V_{12}$ обеспечивает обмен энергией и позволяет системе прийти в состояние теплового равновесия.

В силу малости $V_{12}$ мы можем им пренебречь при записи энергии и считать Hamilton-функцию аддитивной:
\begin{equation}
H(p, q) \approx H_1(p_1, q_1) + H_2(p_2, q_2)
\end{equation}

\subsection{Вывод функции распределения подсистемы}

Наша задача — найти функцию распределения для тела $\rho_1(p_1, q_1)$, зная, что вся составная система $(1+2)$ описывается микроканоническим распределением (МКР). Для полной изолированной системы вероятность найти её в элементе фазового объема $d\Gamma$ равна:
\begin{equation}
dw_{1+2} = \rho_{1+2}(p, q) d\Gamma = \frac{1}{\Omega(E)} \cdot \Delta(H(p, q) - E) \cdot d\Gamma
\end{equation}
Здесь $\Omega(E)$ — статистический вес всей системы, а $\Delta$ — дельта-образная функция, выделяющая слой энергии.

Фазовое пространство полной системы является декартовым произведением пространств каждой подсистемы: $\Gamma = \Gamma_1 \times \Gamma_2$, соответственно $d\Gamma = d\Gamma_1 \cdot d\Gamma_2$.
Чтобы найти вероятность $dw_1$ для нашего тела находиться в определенной точке своего фазового пространства, мы должны проинтегрировать полную вероятность по всем возможным состояниям среды:
\begin{equation}
dw_1(p_1, q_1) = \rho_1(p_1, q_1) d\Gamma_1 = \int_{\Gamma_2} \rho_{1+2}(p_1, q_1, p_2, q_2) d\Gamma_1 d\Gamma_2
\end{equation}

Отсюда следует, что функция распределения подсистемы $\rho_1$ получается интегрированием полной плотности вероятности по переменным среды $(p_2, q_2)$:
\begin{equation}
\rho_1(p_1, q_1) = \frac{1}{\Omega(E)} \int \Delta(H_1(p_1, q_1) + H_2(p_2, q_2) - E) d\Gamma_2
\end{equation}

Обозначим энергию тела в данной точке фазового пространства как $E_1 = H_1(p_1, q_1)$. Тогда под интегралом остается условие на энергию среды: $H_2(p_2, q_2) = E - E_1$. Интеграл от этого условия по фазовому объему среды — это, по определению, статистический вес среды $\Omega_2$ при энергии $E - E_1$:
\begin{equation}
\rho_1(p_1, q_1) = \frac{\Omega_2(E - E_1)}{\Omega(E)}
\end{equation}

\subsection{Разложение по малости энергии подсистемы}

Мы получили отношение статвесов. Теперь воспользуемся условием, что наше тело намного меньше среды по размеру и количеству степеней свободы. Это эквивалентно тому, что энергия тела $E_1$ много меньше полной энергии системы $E$: $E_1 \ll E$.

Напрямую разлагать статвес $\Omega_2$ нельзя, так как это очень быстро растущая функция. Но мы можем разложить её логарифм — энтропию среды $S_2$. 
\begin{equation}
\ln \Omega_2(E - E_1) = S_2(E - E_1) \approx S_2(E) - \left( \frac{\partial S_2(E)}{\partial E} \right) \cdot E_1 + \dots
\end{equation}

Вспоминая термодинамическое определение температуры, мы видим, что производная энтропии среды по её энергии — это обратная температура среды $1/T_2$. Поскольку система находится в равновесии, температура среды совпадает с температурой тела: $T_2 = T$. Получаем:
\begin{equation}
S_2(E - E_1) \approx S_2(E) - \frac{E_1}{T}
\end{equation}

Переходя обратно от логарифма к самой функции, получаем для статвеса среды:
\begin{equation}
\Omega_2(E - E_1) = e^{S_2(E - E_1)} \approx e^{S_2(E)} \cdot e^{-E_1/T} = \Omega_2(E) \cdot e^{-E_1/T}
\end{equation}

Подставляя это в выражение для $\rho_1$, находим:
\begin{equation}
\rho_1(p_1, q_1) = \frac{\Omega_2(E)}{\Omega(E)} \cdot e^{-H_1(p_1, q_1)/T}
\end{equation}

Поскольку множитель перед экспонентой не зависит от микроскопического состояния тела (он определяется только полной энергией и параметрами среды), его можно рассматривать как нормировочную константу $1/Z$. Окончательно, убирая индекс 1, так как переменные среды больше не фигурируют, мы получаем \textbf{распределение Гиббса}:
\begin{equation}
\rho(p, q) = \frac{1}{Z} e^{-H(p, q)/T}
\end{equation}
где $Z$ — статистическая сумма, определяемая из условия нормировки.


% --- SECTION 5: Статистическая сумма и термодинамические потенциалы ---
\section{Статистическая сумма и связь со свободной энергией}

Мы с вами получили фундаментальный результат — каноническое распределение Гиббса. Давайте зафиксируем его и разберем структуру тех величин, которые в него входят. Функция распределения плотности вероятности в фазовом пространстве для системы, находящейся в контакте с термостатом, имеет вид:

\begin{equation}
\rho(p, q) = \frac{1}{Z} e^{-\frac{H(p, q)}{T}}
\end{equation}

Здесь $H(p, q)$ — это Гамильтониан нашей системы, а $T$ — температура, задаваемая внешней средой. Величина $Z$ является нормировочной константой, которую в статистической физике принято называть \textbf{статистической суммой} (или интегралом по состояниям в классическом случае). Она находится из условия нормировки вероятности: сумма (или интеграл) всех вероятностей должна быть равна единице.

\begin{equation}
\int \rho(p, q) d\Gamma = 1 \implies Z = \int e^{-\frac{H(p, q)}{T}} d\Gamma
\end{equation}

Это важнейшее понятие. Если в микроканоническом ансамбле основой основ был статвес $\Omega$, то здесь всё определяется через статистическую сумму $Z$.

\subsection{Энтропия Гиббса}

Теперь возникает вопрос: а совпадает ли эта «новая» статистика с той термодинамикой, к которой мы привыкли? В микроканоническом распределении мы вводили энтропию как логарифм статвеса. Но теперь у нас нет фиксированной энергии, она флуктуирует. Как определить энтропию здесь?

Для этого вводится более общее понятие — \textbf{энтропия Гиббса}. Она определяется как среднее значение от логарифма функции распределения, взятое с обратным знаком:

\begin{equation}
S = - \langle \ln \rho \rangle = - \int \rho(p, q) \ln \rho(p, q) d\Gamma
\end{equation}

Давайте проверим, к чему приведет это определение, если мы подставим в него каноническое распределение. Сначала найдем логарифм плотности вероятности:

\begin{equation}
\ln \rho = \ln \left( \frac{1}{Z} e^{-\frac{H}{T}} \right) = - \ln Z - \frac{H(p, q)}{T}
\end{equation}

Теперь подставим это выражение в интеграл для энтропии (формула 3):

\begin{equation}
S = - \int \rho(p, q) \left( - \ln Z - \frac{H(p, q)}{T} \right) d\Gamma = \ln Z \int \rho d\Gamma + \frac{1}{T} \int H \rho d\Gamma
\end{equation}

Первый интеграл $\int \rho d\Gamma$ в силу условия нормировки равен единице. Второй интеграл $\int H \rho d\Gamma$ — это не что иное, как среднее значение энергии системы, которое мы обозначим как $\langle E \rangle$. В термодинамике мы именно эту величину называем внутренней энергией. Таким образом, мы получаем соотношение:

\begin{equation}
S = \ln Z + \frac{\langle E \rangle}{T}
\end{equation}

\subsection{Свободная энергия Гельмгольца}

Умножим обе части полученного уравнения на температуру $T$:

\begin{equation}
TS = T \ln Z + \langle E \rangle \implies -T \ln Z = \langle E \rangle - TS
\end{equation}

Вспомним классическую термодинамику. Величина $E - TS$ — это термодинамический потенциал, который называется \textbf{свободной энергией} (или энергией Гельмгольца), обозначаемой буквой $F$. Отсюда мы получаем еще одно фундаментальное соответствие между статистикой и термодинамикой:

\begin{equation}
F = -T \ln Z
\end{equation}

Это соотношение является мостиком, который позволяет нам, вычислив статистическую сумму $Z$ через микроскопические параметры системы, сразу найти свободную энергию. А зная свободную энергию как функцию температуры, объема и числа частиц $F(T, V, N)$, мы можем по стандартным правилам термодинамики найти все остальные величины: энтропию, давление, химический потенциал и так далее.

\subsection{Связь распределения со свободной энергией}

Используя полученную связь $F = -T \ln Z$, мы можем переписать константу $1/Z$ в распределении Гиббса. Очевидно, что $1/Z = e^{F/T}$. Тогда плотность вероятности принимает вид:

\begin{equation}
\rho(p, q) = e^{\frac{F - H(p, q)}{T}}
\end{equation}

Эта запись очень удобна и часто используется в теоретических выкладках. Посмотрите, как изящно всё получилось: вся сложность взаимодействия системы с огромным термостатом, все эти бесчисленные степени свободы среды «схлопнулись» и вошли в наш расчет всего через два параметра — температуру $T$ и свободную энергию $F$ (которая сама является функцией $T$). Термостат задает температуру и формирует именно такой экспоненциальный вид распределения.

В этом и заключается мощь метода Гиббса: нам больше не нужно следить за средой, нам достаточно знать, что она поддерживает систему в равновесии при заданной температуре.

% --- SECTION 6: Практическое применение канонического ансамбля ---
\section{Пример: Классический идеальный газ Больцмана}

Теперь давайте проиллюстрируем всё вышесказанное конкретным вычислением. Самый поучительный и важный пример для понимания того, как работает формализм канонического распределения, — это классический идеальный газ. 

Гамильтониан такой системы представляет собой просто сумму кинетических энергий всех частиц, так как потенциальной энергией взаимодействия в идеальном газе мы пренебрегаем:
\begin{equation}
H(p, q) = \sum_{i=1}^{N} \frac{\vec{p}_i^{\,2}}{2m}
\end{equation}
Здесь $\vec{p}_i$ — импульс $i$-й частицы. Как мы видим, Гамильтониан явно не зависит от координат частиц $\vec{r}_i$.

\subsection{Вычисление статистической суммы}

Наша цель — посчитать статистическую сумму $Z$. Согласно определению канонического ансамбля:
\begin{equation}
Z = \int e^{-\frac{H(p, q)}{T}} d\Gamma
\end{equation}
Вспомним выражение для элемента фазового объема $d\Gamma$ для системы из $N$ тождественных частиц. С учетом квантовой нормировки и множителя Гиббса для неразличимости частиц, имеем:
\begin{equation}
d\Gamma = \frac{d^3p_1 \dots d^3p_N \cdot d^3r_1 \dots d^3r_N}{N! (2\pi\hbar)^{3N}}
\end{equation}
Подставляя Гамильтониан (1) в экспоненту, мы замечаем важное свойство: экспонента от суммы превращается в произведение экспонент. Это позволяет нам разбить полный интеграл на произведение независимых интегралов по каждой частице.

Сначала разберемся с пространственной частью. Поскольку Гамильтониан не зависит от координат $\vec{r}_i$, каждый интеграл $\int d^3r_i$ по объему ящика дает просто объем $V$. Для $N$ частиц это даст множитель $V^N$.

Теперь перейдем к импульсной части. Статсумма распадается в произведение $N$ одинаковых интегралов по импульсам каждой частицы:
\begin{equation}
Z = \frac{V^N}{N! (2\pi\hbar)^{3N}} \left( \int e^{-\frac{p^2}{2mT}} d^3p \right)^N
\end{equation}

Интеграл по импульсу одной частицы удобно вычислять в сферических координатах ($d^3p = 4\pi p^2 dp$):
\begin{equation}
I_p = \int_0^\infty 4\pi p^2 e^{-\frac{p^2}{2mT}} dp
\end{equation}
Это стандартный гауссов интеграл. Напомним общую формулу:
\begin{equation}
\int_0^\infty e^{-\alpha x^2} dx = \frac{1}{2} \sqrt{\frac{\pi}{\alpha}}, \quad \int_0^\infty x^2 e^{-\alpha x^2} dx = \frac{1}{4\alpha} \sqrt{\frac{\pi}{\alpha}}
\end{equation}
В нашем случае параметр $\alpha = \frac{1}{2mT}$. Применяя формулу, получаем:
\begin{equation}
I_p = 4\pi \cdot \frac{2mT}{4} \sqrt{2\pi mT} = (2\pi mT)^{3/2}
\end{equation}

\subsection{Тепловая длина волны и окончательный вид статсуммы}

Подставим полученный интеграл в выражение для $Z$:
\begin{equation}
Z = \frac{V^N}{N!} \left( \frac{(2\pi mT)^{3/2}}{(2\pi\hbar)^3} \right)^N = \frac{V^N}{N!} \left( \frac{2\pi mT}{h^2} \right)^{3N/2}
\end{equation}
Здесь мы перешли от перечеркнутой постоянной Планка $\hbar$ к обычной $h = 2\pi\hbar$. Посмотрите на величину в скобках. Она имеет размерность обратного объема. Мы можем снова ввести тепловую длину волны де Бройля:
\begin{equation}
\lambda_T = \frac{h}{\sqrt{2\pi mT}}
\end{equation}

% Место под рисунок
% Комментарий: Лектор схематично изображает волновой пакет частицы в ящике, поясняя смысл \lambda_T как пространственного масштаба квантовой неопределенности.
% Таймкод: [0:59:50]

Тогда статистическая сумма идеального классического газа принимает исключительно компактный и красивый вид:
\begin{equation}
Z = \frac{V^N}{N! \lambda_T^{3N}} = \frac{1}{N!} \left( \frac{V}{V_Q} \right)^N
\end{equation}
где $V_Q = \lambda_T^3$ — уже знакомый нам квантовый объем. Используя формулу Стирлинга $\ln N! \approx N \ln(N/e)$, мы можем записать:
\begin{equation}
Z \approx \left( \frac{eV}{N V_Q} \right)^N
\end{equation}

\subsection{Сравнение микроканонического и канонического подходов}

Давайте сравним, что мы сделали сейчас и что делали в начале лекции. 
В микроканоническом ансамбле нам приходилось считать \textit{объем многомерного шара} — это довольно громоздкая геометрическая задача, требующая знания специальных функций (Гамма-функции) и аккуратного обращения с оболочками в фазовом пространстве.

В каноническом же распределении вычисление свелось к \textit{произведению одномерных гауссовых интегралов}. Математически это гораздо проще. При этом результат для свободной энергии $F = -T \ln Z$ будет полностью эквивалентен тому, что мы получали через энтропию Больцмана, если перейти к пределу больших систем.

Более того, здесь мы сразу видим физические условия применимости. Наш расчет Гамильтониана как суммы кинетических энергий справедлив, пока газ разрежен. Условие классичности, как мы уже отмечали, записывается через отношение удельного объема к квантовому:
\begin{equation}
v = \frac{V}{N} \gg V_Q
\end{equation}
Если температура стремится к нулю, $V_Q$ растет как $T^{-3/2}$ и в какой-то момент станет соизмерим с $v$. Это температура вырождения $T_0$, при которой газ становится квантовым. Но пока мы находимся в режиме $T \gg T_0$, каноническое распределение Гиббса дает нам максимально простой и эффективный способ описания системы.

На следующей лекции мы, пользуясь этой функцией $F(T, V, N)$, выведем все термодинамические соотношения, включая формулу Саккура-Тетроде для энтропии, и убедимся в их корректности. На сегодня всё.

\end{document}
